% Irodov Part IV §4.2 original 126--137. Source PDF pp.181--182,
% printed pp.172--173. Answers PDF pp.319--320.
\electricoscproblem{126}{172}
电感 $L=0.36\,\mathrm H$、有功电阻
$R=40\,\Omega$ 的线圈与电容 $C$ 串联，
外加正弦电压角频率 $\omega=314\,\mathrm{s^{-1}}$、
幅度 $V_m=180$V。
调节 $C$ 使\emph{整个线圈}端电压幅度最大；
求最佳电容以及对应线圈、电容电压幅度。
\exsolution
线圈幅度阻抗
$|Z_L|=\sqrt{R^2+(\omega L)^2}$
在调节 $C$ 时不变；线圈端电压
$V_L=V_m|Z_L|/
\sqrt{R^2+(\omega L-1/(\omega C))^2}$。
分母在总反应抗抵消时最小：
$1/(\omega C)=\omega L$，
故
\[
\boxed{C_*=\frac1{\omega^2L}
\simeq28\,\mathrm{\mu F}}.
\]
此时 $I_m=V_m/R$，
\[
\boxed{V_{{\rm coil},m}=
V_m\sqrt{1+\left(\frac{\omega L}{R}\right)^2}
\simeq0.54\,\mathrm{kV}},
\quad
\boxed{V_{C,m}=\frac{V_m\omega L}{R}
\simeq0.51\,\mathrm{kV}}.
\]
\excheck 两器件端电压的幅度均可高于源幅度，源电压为它们的相量和。

\electricoscproblem{127}{172}
介质有微弱导电性的电容器，电容为 $C$、漏电电阻为 $R$，
两端接正弦电压
$V(t)=V_m\cos\omega t$。
导线电阻忽略，求稳态供电导线中的电流随时间的变化。
\exsolution
漏电传导电流
$I_R=V/R=(V_m/R)\cos\omega t$；
位移电流
$I_C=C\dot V=-\omega CV_m\sin\omega t$。
总供电电流
\[
I(t)=\frac{V_m}{R}\cos\omega t-\omega CV_m\sin\omega t.
\]
写成相位形式
\[
\boxed{I(t)=I_m\cos(\omega t+\varphi)},\qquad
\boxed{I_m=\frac{V_m}{R}\sqrt{1+(\omega RC)^2}},
\quad
\boxed{\tan\varphi=\omega RC}.
\]
\excheck 导线电流包括漏电传导与充放电电流；
它超前电压，且 $\omega\to0$ 回到纯电阻漏电流。

\electricoscproblem{128}{172}
LC 回路电容为 $C$、自感为 $L_1$；
其线圈与另一只理想短路线圈（自感为 $L_2$）互感为 $M$。
短路线圈有功电阻忽略。
求该 LC 回路的固有角频率。
\exsolution
两线圈的有向磁链
$\Psi_1=L_1I_1+MI_2$、
$\Psi_2=MI_1+L_2I_2$。
次级线圈短路且无电阻，
其磁链变化为零；若初始磁链为零，
$\Psi_2=0$，所以
$I_2=-(M/L_2)I_1$。
代回初级磁链
$\Psi_1=(L_1-M^2/L_2)I_1$，
等效自感
$L_{\rm eff}=L_1-M^2/L_2$。
故
\[
\boxed{\omega_0=
\frac1{\sqrt{C(L_1-M^2/L_2)}}
=\sqrt{\frac{L_2}{C(L_1L_2-M^2)}}}.
\]
\excheck 由磁能正定性 $M^2<L_1L_2$，
等效电感必须为正；电磁耦合缩小初级电感，从而提高频率。

\electricoscproblem{129}{172}
串联谐振回路接交流电源，在电流共振时
电容器两端电压幅度是外加电源幅度的 $n$ 倍。
求品质因数 $Q$。
\exsolution
串联电流共振频率
$\omega_0=1/\sqrt{LC}$，
电路阻抗最小为 $R$，
$I_m=V_m/R$。
电容电压幅度与源电压幅度之比
\[
n=\frac{I_m/(\omega_0 C)}{V_m}
=\frac1{\omega_0CR}
=\frac{\omega_0L}{R}.
\]
若品质因数按\emph{实际阻尼频率}
$\omega_d=\sqrt{\omega_0^2-\beta^2}$ 定义：
$Q=\omega_d/(2\beta)$，
$\beta=R/(2L)$，
则
\[
\boxed{Q=\sqrt{\left(\frac{\omega_0L}{R}\right)^2
-\frac14}
=\sqrt{n^2-\frac14}}.
\]
\excheck 弱阻尼时 $Q\simeq n$；区别来自采用 $\omega_d$ 还是 $\omega_0$ 定义 Q，而不是测量电压出现错误。

\electricoscproblem{130}{172}
串联 RLC 回路的电源角频率与电压幅度保持不变，
调整自感 $L$ 使电流达到最大幅度。
若此时自感增加到原来的 $n$ 倍，电流幅度降为原来的 $1/\eta$，
求原共振状态的品质因数。
\exsolution
调节 $L$ 使
$\omega L=1/(\omega C)$，
原最大电流 $I_m=V_m/R$。
自感换为 $nL$ 后，
反应抗
$X'=\omega(nL)-1/(\omega C)
=\omega L(n-1)$，
所以
\[
\eta^2=
\frac{R^2+[\omega L(n-1)]^2}{R^2}
=1+\left(\frac{\omega L}{R}\right)^2(n-1)^2.
\]
由第129题的品质因数约定
$Q^2=(\omega L/R)^2-1/4$，
得
\[
\boxed{Q=
\sqrt{\frac{\eta^2-1}{(n-1)^2}-\frac14}}.
\]
\excheck 需有 $n\ne1$、$\eta\ge1$；
此关系还给出数据相容性 $Q^2\ge0$。

\electricoscproblem{131}{172}
串联 RLC 回路交流源电压幅度恒定，
改变电源频率时，在
$\omega_1,\omega_2$ 两处电流幅度都比电流共振峰小 $n$ 倍。
求：(a) 共振角频率；(b) 品质因数。
\exsolution
电流幅度
$I(\omega)=V_m/\sqrt{R^2+(\omega L-1/(\omega C))^2}$。
在电流峰 $\omega_0^2=1/(LC)$，
$I_{\max}=V_m/R$。
半峰倍率条件
\[
\left(\frac{\omega_0^2-\omega_i^2}{\omega_i}\right)^2
=(n^2-1)\left(\frac RL\right)^2.
\]
较低和较高两正根对应相反符号，
解得
$\omega_1\omega_2=\omega_0^2$ 与
$\omega_2-\omega_1=(R/L)\sqrt{n^2-1}$。
故
\[
\boxed{\omega_0=\sqrt{\omega_1\omega_2}}.
\]
按 $Q=\omega_dL/R$，$\omega_d^2=\omega_0^2-(R/2L)^2$：
\[
\boxed{Q=
\sqrt{\frac{\omega_1\omega_2(n^2-1)}
{(\omega_2-\omega_1)^2}-\frac14}}.
\]
\excheck 这里 $\omega_1,\omega_2$ 必须位于峰的两侧且 $n>1$。

\electricoscproblem{132}{172}
证明弱阻尼串联谐振回路的品质因数
$Q\simeq\omega_0/\Delta\omega$，
其中 $\Delta\omega$ 是电流谐振曲线在
\emph{峰值幅度的 $1/\sqrt2$（半功率）}处的频宽。
\exsolution
在 $\omega$ 接近 $\omega_0$ 时，
串联反应抗
$X(\omega)=\omega L-1/(\omega C)
\simeq2L(\omega-\omega_0)$。
所以
\[
\frac{I(\omega)^2}{I_{\max}^2}
\simeq\frac{R^2}{R^2+
4L^2(\omega-\omega_0)^2}.
\]
半功率点
$I^2=I_{\max}^2/2$
对应
$|2L(\omega-\omega_0)|=R$，
两侧点间宽度
$\Delta\omega=R/L$。
弱阻尼
$Q\simeq\omega_0L/R$，
故
\[
\boxed{Q\simeq\frac{\omega_0}{\Delta\omega}}.
\]
\excheck 底本文字抽取把“$\sqrt2$ times less”中的根号漏掉，只剩“2 times less”；若按幅度真正降为1/2，则频宽系数不同，不能推导此式。

\electricoscproblem{133}{172--173}
串联 RLC 回路先后接入电压幅度相同、频率分别为
$\omega_0$ 和 $\eta\omega_0$ 的两个交流源，
其中 $\omega_0=1/\sqrt{LC}$。
若品质因数为 $Q$，
求两情况下电流幅度之比 $I_0/I$；
并在 $\eta=1.10$ 时分别算
$Q=10$ 和 $Q=100$。
\exsolution
在 $\omega_0$ 处阻抗是 $R$，
$I_0=V_m/R$。
在 $\eta\omega_0$ 处的反应抗
\[
X=\eta\omega_0L-\frac1{\eta\omega_0C}
=\omega_0L\frac{\eta^2-1}{\eta}.
\]
故
\[
\frac{I_0}{I}=
\sqrt{1+\left(\frac X R\right)^2}.
\]
用本书阻尼品质因数定义
$\omega_0^2L^2/R^2=Q^2+1/4$，
得
\[
\boxed{\frac{I_0}{I}=
\sqrt{1+\frac{(Q^2+1/4)(\eta^2-1)^2}{\eta^2}}}.
\]
代 $\eta=1.10$：
\[
\boxed{I_0/I\simeq2.2\quad(Q=10)},\qquad
\boxed{I_0/I\simeq19\quad(Q=100)}.
\]
\excheck 高品质因数回路对较小的失谐更敏感；
两个数值均符合书末。

\electricoscproblem{134}{173}
蓄电池以恒定直流电流 $I_0$ 充电需 $t_0$ 小时。
若改用半波整流交流电流充电，
整流电流的\emph{有效值}仍为 $I_0$，
问需要多长时间，忽略实际电池充电效率的变化。
\exsolution
半波整流正弦电流一个周期中
$I(t)=I_m\sin\omega t$（$0<\omega t<\pi$），
另一半周期为零。
有效值平方
$I_{\rm rms}^2=(2\pi)^{-1}\int_0^\pi I_m^2\sin^2u\,\mathrm du=I_m^2/4$，
故 $I_m=2I_0$。
周期平均值
$\overline I=(2\pi)^{-1}\int_0^\pi I_m\sin u\,\mathrm du=I_m/\pi=2I_0/\pi$。
相同总电荷
$I_0t_0=\overline It$，所以
\[
\boxed{t=\frac\pi2t_0}.
\]
\excheck 有效电流相同不等于平均输送电荷速率相同；半波电流的峰值与平均值都必须单独计算。

\electricoscproblem{135}{173}
一股周期电流在整个周期的平均值为 $I_0$。
求电流的有效值：
(a) 如原图4.34，电流在每一周期从零\emph{线性增加到峰值}后瞬间归零的锯齿波；
(b) $I(t)\propto|\sin\omega t|$。
\exsolution
(a) 令周期为 $T$，
$I(t)=I_m t/T$（$0<t<T$）。
平均
$I_0=T^{-1}\int_0^T I_m t/T\,\mathrm dt=I_m/2$；
有效值
\[
I_{\rm rms}^2=\frac1T\int_0^T
\left(\frac{I_mt}{T}\right)^2\mathrm dt
=\frac{I_m^2}{3}.
\]
因 $I_m=2I_0$，
\[
\boxed{I_{{\rm rms},a}=\frac2{\sqrt3}I_0\simeq1.15I_0}.
\]
(b) 对全波整流正弦
$\overline I=(2/\pi)I_m=I_0$，
$I_{\rm rms}=I_m/\sqrt2$，
故
\[
\boxed{I_{{\rm rms},b}=\frac{\pi}{2\sqrt2}I_0
\simeq1.11I_0}.
\]
\excheck (a) 的波形来自原图4.34，不能误认成三角上下对称波而改变积分。

\electricoscproblem{136}{173}
线圈 $L=7\,\mathrm{mH}$、
$R=44\,\Omega$，先接直流源 $V_0$，
再接有效电压同为 $V=V_0$ 的正弦交流源。
求交流频率 $\nu$，
使线圈交流有功热功率只有直流情形的 $1/\eta$，
其中 $\eta=5.0$。
\exsolution
直流稳定时 $P_{\rm DC}=V_0^2/R$；
交流时总阻抗
$Z=\sqrt{R^2+(2\pi\nu L)^2}$，
有功功率
$P_{\rm AC}=V_0^2R/Z^2$。
比值
$P_{\rm DC}/P_{\rm AC}
=Z^2/R^2=1+(2\pi\nu L/R)^2=\eta$，
所以
\[
\boxed{\nu=\frac{R}{2\pi L}\sqrt{\eta-1}
=\frac{44}{2\pi(0.007)}(2)
\simeq2.0\,\mathrm{kHz}}.
\]
\excheck 电压给的是有效值，因而功率表达式不用再额外除2。

\electricoscproblem{137}{173}
线圈感抗为 $X_L=30\,\Omega$，
总阻抗模为 $Z=50\,\Omega$，
接有效电压 $V=100$V 的交流电源。
求电流相对电压的相位差，以及线圈平均发热功率。
\exsolution
线圈阻抗为 $R+iX_L$，
$R=\sqrt{Z^2-X_L^2}
=\sqrt{2500-900}=40\,\Omega$。
相位滞后角
\[
\boxed{\varphi=\arctan\frac{X_L}{R}
=\arctan\frac{30}{40}
\simeq37^\circ}.
\]
电流有效值 $I=V/Z=2.0$A，
热功率
\[
\boxed{P=I^2R=(2.0)^2(40)
=160\,\mathrm W=0.16\,\mathrm{kW}}.
\]
\excheck 感性线圈中电流\emph{滞后}电压，只有其有功电阻产生焦耳热。
